% --------------------------------------------------
% 2. THEORETICAL FRAMEWORK AND LITERATURE
% --------------------------------------------------
\section{Theoretical Framework and Literature Review}
\label{sec:theory}

\subsection{The Economics of Rider Abandonment}

A Grand Tour lasts three weeks and consists of 21 stages. At the start of each stage, every rider still in the race faces a decision: continue or abandon. This paper frames that decision as a rational choice under uncertainty, where the rider and his team compare the expected benefit of continuing (GC position gains, UCI points, sponsor exposure, contract value) against the expected cost (fatigue, injury risk, reduced future performance).

This framing builds on tournament theory as developed by \textcite{rosen1981superstars} and \textcite{lazear1981rank}. Prize structures in tournaments are typically convex: the reward gap between first and second place far exceeds the gap between 50th and 51st. In Grand Tour cycling, this means riders near the top of the GC standings have disproportionately more to gain from staying in the race, while a rider in 80th place may rationally withdraw and save energy for future events. \textcite{szymanski2003economicdesign} provides a broader framework for understanding how the design of sporting contests shapes competitive incentives, and the same logic applies to the stage-by-stage structure of Grand Tours.

A distinctive feature of Grand Tours, compared to the single-contest tournaments in \textcite{lazear1981rank}, is their \emph{nested} or \emph{multi-stage} structure. Each individual stage functions as a tournament (with stage wins and time bonuses as prizes), the 21-stage GC classification is a tournament across stages, and the season-long UCI ranking is a tournament across races. These layers create competing incentive gradients. A domestique may have no stake in the GC tournament but a strong interest in winning an individual stage; a GC contender cares little about individual stage prizes but enormously about the overall classification. This means the same physical event (a mountain stage, for example) generates different economic incentives for different riders depending on which layer of the tournament structure they are competing in.

The multi-stage structure also introduces an \emph{option value} dimension to the continuation decision. At each stage, a rider who is still in the race holds an option: the option to win a future stage, to move up in the GC standings, or to be available for a team strategic move later in the race. Continuing preserves this option; abandoning extinguishes it permanently. The value of the option depends on how many stages remain (higher early in the race), the rider's current position (higher for riders still in GC contention), and the expected difficulty of upcoming stages (lower if brutal mountains lie ahead). This option-value logic predicts that abandonment hazard should increase over the course of a race (as fewer stages remain, the option value of continuing diminishes) and should be higher on mountain stages (which impose a large exercise cost on the option). Both predictions are broadly consistent with the descriptive patterns in the data.

This tournament logic extends to the team level. As discussed in the course lectures on sport clubs as profit-maximising firms (Lecture F2), teams operate with budget constraints and strategic objectives \parencite{blair2012sports}. A team may withdraw a domestique who has fulfilled his support role and is carrying an injury, preserving his physical capital for races where his marginal product will be higher. Conversely, if a GC contender abandons, the domestiques allocated to protect him lose their primary function. The team thus acts as a firm making resource-allocation decisions about its riders, weighing each rider's expected marginal contribution in the current race against his expected contribution in future events.

\subsection{Labour Markets and Contracts in Professional Cycling}

Professional cyclists operate in a labour market that shares features with other professional sports labour markets. \textcite{kahn2000labormarket} argues that professional sports provide a uniquely useful laboratory for studying labour-market questions because of the availability of detailed performance data and the repeated nature of competition. Grand Tour cycling fits this framework: riders produce observable output (stage results, GC rankings, race completions) that signals their quality to teams negotiating future contracts. \textcite{rottenberg1956baseball} established that even in highly regulated sports labour markets, players respond to economic incentives in predictable ways, and the same logic applies to cyclists deciding whether to continue or withdraw.

Within this framework, abandoning a race is a negative performance signal. A rider who does not finish provides less information to the market about his ability, and the act of abandoning may be interpreted as evidence of fragility or insufficient preparation. However, the reputational cost of a DNF depends on context: a domestique who abandons after serving his team leader for two weeks faces different labour-market consequences than a GC contender who abandons while in fifth place (Lectures F4 and F5).

The contract structure in professional cycling adds a further dimension. \textcite{krautmann2002contracts} show that in Major League Baseball, contract length reflects expected future performance and the information environment. In cycling, riders on short-term contracts approaching renewal face different incentives from those on long-term deals. A rider in a contract year may push through pain and fatigue to signal durability, knowing that a DNF could weaken his bargaining position. Conversely, a rider with a secure multi-year contract may abandon more freely, since the immediate reputational cost is buffered by the existing agreement. \textcite{kahn1993freeagency} provides evidence from baseball that the structure of labour-market institutions (free agency, reserve clauses) affects player behaviour, and analogous institutional features in cycling (the transfer system, minimum wage regulations, team licence obligations) shape the incentive environment in which abandonment decisions are made.

The dataset includes each rider's historical DNF rate. This variable can be interpreted in two ways: a high rate may indicate a ``fragile type'' who faces higher continuation costs, or it may reflect that a rider with several past DNFs faces lower marginal reputational cost from one more, because the signal has already been sent. Both interpretations predict a positive association with current abandonment hazard.

Race prestige provides another source of variation. The Tour de France attracts the largest television audiences, the most media coverage, and the most sponsor attention, so the marginal revenue product of continued participation is likely higher in the Tour than in the Giro or Vuelta \parencite{blair2012sports}. This connects to the course material on advertisement and TV rights (Lecture F3): broadcasting revenues and sponsor contracts are tied to participation and visibility, giving riders and teams stronger incentives to persist in higher-prestige events. \textcite{mignot2016cycling} documents the Tour's dominant commercial position among Grand Tours, providing historical context for this incentive asymmetry.

\subsection{Related Empirical Work}

Several studies in sports economics have used hazard or duration models to analyse decisions with a structure similar to rider abandonment. The course material on managers, coaches, and team performance (Lecture F6) discusses how managerial tenure in football can be modelled as a survival process, where dismissal probability depends on match results, league position, and club characteristics. The underlying logic is the same as in this paper: an event (dismissal or abandonment) occurs at some point during a spell (a season or a race), and covariates shift the probability of that event. \textcite{boeri2008italianjob} study career concerns among referees in Italian football and show how incentive structures lead to strategic behaviour that can be inferred from patterns in the data, an approach conceptually similar to what this paper does with rider-level proxies. \textcite{rosen2001labourmarkets} provide a broader treatment of how professional sports labour markets generate observable signals that researchers can exploit.

The survival-analysis methodology used in this paper has parallels in other sports contexts. Career-duration studies in American football, baseball, and basketball use hazard models to estimate the determinants of career length, treating retirement or release as the event of interest. The distinguishing feature of the present application is that the ``spell'' is a single race (21 stages) rather than an entire career, which gives much finer temporal resolution and allows stage-level covariates that vary within the spell. \textcite{courty2003ticketresale} demonstrates from a different angle that outcomes within sporting events carry economic value: his work on ticket resale shows that spectators pay a premium for uncertainty of outcome, and rider abandonment resolves that uncertainty, reducing the entertainment and therefore commercial value of remaining stages.

Outside the course material, there is a growing literature on cycling economics \parencite{torgler2007cycling}, but most of it focuses on doping, team strategy, or stage-win determinants. Very few papers model within-race attrition from an explicitly economic perspective. This is surprising, because Grand Tours generate rich stage-by-stage panel data well suited to survival analysis, and because abandonment has clear economic consequences: it reshuffles GC standings, redistributes prize money and UCI points, and affects the information content of race results as signals of rider quality. This paper aims to fill that gap by applying a standard econometric framework to a novel dataset, and by interpreting the results through the lens of tournament incentives and labour-market signalling.
